% Working continuation toward A.4.57.  These finite lemmas are proved;
% the complete CM compiler receives its result wrapper after review.
\subsection{Finite cone and peripheral descent for the branch transport}

Retain the actual orbifold
$\mathcal E=C_2(b_1)*C_2(b_2)*C_2(b_3)*
F(b_0,t_1,t_2,t_3)$, its original index-two surface subgroup
$K=\Pi_{1,9}$, and the different quotient
$\mathcal E\to\Omega=C_2(t_1)*\cdots*C_2(t_4)$
of (FBO.7)--(FBO.11).

\par\medskip\noindent\textbf{An explicit peripheral square-root detector.}\quad
Let $q_0:\mathcal E\twoheadrightarrow D_0$ be any specified
finite ordinary quotient, and set
$r=\operatorname{ord}_{D_0}(q_0(b_0))$.
Use the actual free-generator character
$e_0:\mathcal E\to\mathbb Z$ sending $b_0$ to one and the
other six free-product generators to zero.  Form the actual
joint image and its homology refinement
\[
 \begin{gathered}
 D=\operatorname{im}(q_0,e_0\bmod2r),\qquad
 L=\ker(\mathcal E\to D),\qquad
 D^+=\mathcal E/(L'L^3),\qquad D^+\twoheadrightarrow D
                                      \twoheadrightarrow D_0.
 \end{gathered}                                                   \tag{OBT.1}
\]
For every $\lambda\in\widehat{\mathbb Z}^{\times}$ and every
$v$ in the specified finite group $D^+$ with
$v^2=b_0^{2\lambda}$, the image of $v$ in $D_0$ is exactly
$q_0(b_0)^\lambda$.  More generally the same statement holds
with the actual prescribed conjugator on both sides of that
square equation.  Also
\[
 \operatorname{im}\bigl(C_{D^+}(b_0^2)\to D\bigr)
            =\langle q(b_0)\rangle.
                                                               \tag{OBT.2}
\]
If $d_0=|D_0|$, the construction has the explicit bound
\[
             |D^+|\leq 2d_0^2\,3^{1+12d_0^2}.
                                                               \tag{OBT.3}
\]


\begin{proof}
The order of $q(b_0)$ in the joint image $D$ is exactly $2r$.
Use the graph of groups in the proof of (FBO.17).  A cone
vertex in the graph for $L$ has stabilizer either trivial or
$C_2$, and its positive homology over $\mathbb F_3$ vanishes.
Thus $H_1(L;\mathbb F_3)$ is again the first homology of the
underlying graph.  The graph has $7|D|$ edges and at least
$|D|$ vertices, so its first Betti number is at most
$1+6|D|$.  Since $|D|\leq2r d_0\leq2d_0^2$, this proves
the bound (OBT.3).

The loop $(b_0^2)^r=b_0^{2r}$ is the primitive $b_0$ cycle
of the regular graph at the identity vertex.  Its edge chain
has coefficient one on each of the $2r$ distinct oriented
$b_0$ edges based at $\langle q(b_0)\rangle$.
Its stabilizer under left deck translation is exactly that
cyclic subgroup.  A centralizer element of $b_0^2$ fixes
this homology class.  Conversely each power of $q(b_0)$ has
the centralizing lift given by the same power of $b_0$.
This proves (OBT.2), including $r=1$.

If $v^2=b_0^{2\lambda}$, then $v$ commutes with
$b_0^{2\lambda}$.  Powering by the inverse of $\lambda$
modulo the finite order of $b_0^2$ shows that it commutes
with $b_0^2$.  Thus (OBT.2) gives $v_D=q(b_0)^a$ for an
integer $a$.  The retained $C_{2r}$ coordinate of the square
equation gives $2a=2\lambda\pmod {2r}$, hence
$a=\lambda\pmod r$.  Projection to $D_0$ is exactly the
asserted value.  Conjugating the original square equation
by its named conjugator gives the stated conjugate version.
No division by two in a finite even-order cyclic group has
been used: the doubled exponent coordinate is retained
before passing to the original target.
\end{proof}

\par\medskip\noindent\textbf{The simultaneous cone input.}\quad
Let $D_0$ now be a specified finite quotient of $\Omega$,
viewed through $\mathcal E\to\Omega$.
Join it with the full elementary abelian two-quotient of
$\mathcal E$, whose basis is
$(b_1,b_2,b_3,b_0,t_1,t_2,t_3)$, and denote that actual
image by $D_{\mathrm{cone}}$.
Its kernel is free because every cone generator retains
order two.  Apply the two binary homology covers
(CONE.1)--(CONE.7).  The resulting actual quotient
$T_{\mathrm{cone}}$ has the property that every order-two
element projects into a conjugate of one of the three
cone groups of $D_{\mathrm{cone}}$.
If $q=|D_0|$, the explicit bounds are
\[
 \begin{gathered}
 d_c\leq128q,\qquad r_0=1+9d_c/2,\qquad
 r_1=1+2^{r_0}(r_0-1),\qquad
 |T_{\mathrm{cone}}|=d_c2^{r_0+r_1}.
 \end{gathered}                                                   \tag{OBT.4}
\]
The source has not been replaced by a product of unrelated
cone residues.  The joint character and both homology covers
are taken on the same actual group $\mathcal E$.

Choose $m$ at least the orders of $T_{\mathrm{cone}}$ and
the square-root detector (OBT.1) over $D_0$, and put
\[
                 P=\mathcal E/\kappa_m(\mathcal E).
                                                               \tag{OBT.5}
\]
It maps to both actual finite targets by their regular
permutation representations.  Every automorphism of $P$
therefore sends the three cone generators to elements whose
images in $D_{\mathrm{cone}}$ are conjugates of cone
generators.  The induced automorphism of
$P/[P,P]P^2=H_1(\mathcal E;\mathbb F_2)$ preserves linear
independence.  The images of the three original cone basis
vectors consequently have three different cone types.
They determine one common permutation in $\mathfrak S_3$.
In particular all three transported cone generators die
under the specified projection to $D_0$.

\par\medskip\noindent\textbf{Descent to a characteristic four-cone quotient.}\quad
Let
$Q_t=\Omega/\kappa_t(\Omega)$, for $t\geq2$, and use
$D_0=Q_t$ in (OBT.4)--(OBT.5).
Let $\gamma\in\operatorname{Aut}(P)$ be an actual finite
automorphism with the following specified peripheral data:
\[
 \gamma(b_0^2)=b_0^{2\lambda},\qquad
 \gamma(t_j)\sim t_j^\lambda\quad(1\leq j\leq4)
                           \quad\hbox{in }P,
                                                               \tag{OBT.6}
\]
where each conjugacy is in that same finite group and
$\lambda\in\widehat{\mathbb Z}^{\times}$.
Then $\gamma$ descends under the actual quotient $P\to Q_t$
to an automorphism $\bar\gamma$ of $Q_t$.  It satisfies
\[
 \bar\gamma(b_0)=b_0^\lambda,\qquad
 \bar\gamma(t_j)\sim t_j^\lambda,
 \qquad\chi_T\bar\gamma=\chi_T.
                                                               \tag{OBT.7}
\]
For any $s$ such that
$\kappa_t(\Omega)\subseteq\kappa_s(K_*)$, its restriction
to the even subgroup descends to an actual automorphism of
$W_s=K_*/\kappa_s(K_*)$, preserving both ordered pole
classes with power $\lambda$.


\begin{proof}
The common cone permutation established after (OBT.5) kills
the three cone generators under the composite
$\mathcal E\to P\xrightarrow{\gamma}P\to Q_t$.
Its four root generators are conjugates of the $\lambda$th
powers of involutions, so their squares are one.  The
presentation (FBO.7)--(FBO.9) therefore gives an actual
factorization of this composite through a homomorphism
$\Omega\to Q_t$.

The latter map kills $\kappa_t(\Omega)$.
To prove this exact assertion, use the defining diagonal
embedding of $Q_t$ in the product of copies of
$\mathfrak S_t$, indexed by all homomorphisms
$\Omega\to\mathfrak S_t$.
Composing the new homomorphism with each coordinate gives
another homomorphism $\Omega\to\mathfrak S_t$ and hence
kills $\kappa_t(\Omega)$.  The diagonal embedding is
injective, so the new map kills it as well.
We have therefore proved the quotient square, rather than
assuming that the target has order at most $t$.
The resulting endomorphism of $Q_t$ is surjective because
$P\xrightarrow{\gamma}P\to Q_t$ is surjective.  It is
an automorphism of this finite group.

The first identity in (OBT.6), projected to (OBT.1), gives
the first identity in (OBT.7) by the square-root lemma.
The other identities follow from their literal generator
values.  In particular $\lambda$ is odd, so every $t_j$
still has $\chi_T$ value one.  Thus the even subgroup of
$Q_t$ is preserved.

Write $N_t=\kappa_t(\Omega)$.
Under the specified containment, every homomorphism
$K_*\to\mathfrak S_s$ factors through $K_*/N_t$,
and conversely.  Hence
\[
 \kappa_s(K_*/N_t)=\kappa_s(K_*)/N_t.
                                                               \tag{OBT.8}
\]
This is a characteristic subgroup of the actual even
subgroup.  It proves descent to $W_s$.
The first pole is represented by $b_0$; its image is the
prescribed power in (OBT.7).
The other is represented, up to the explicit even conjugator
in (FBO.11), by $t_1b_0t_1$.
Since $\bar\gamma(t_1)$ is odd, write
$\bar\gamma(t_1)=h t_1$ with $h$ even.
Then the other pole maps to
$h(t_1b_0t_1)^\lambda h^{-1}$.
This retains the second ordered class as well.
\end{proof}

\par\medskip\noindent\textbf{How the hypotheses are supplied by the actual surface input.}\quad
The automorphism in (OBT.6) is obtained from the characteristic
rank-ten surface refinement in (FBO.19)--(FBO.25), applied
with this $m$.  The finite nine-point tower transport that
fixes the named involution normalizes its deck group in
the actual outer surface representation.  The intrinsic
inner-action reconstruction then gives $\gamma$.
Choose the based surface representative by conjugating its
image of the named origin loop $b_0^2$ to the literal
$\lambda$th power.  This is a conjugation inside the
characteristic surface quotient, so it leaves the outer
mapping-class covariance unchanged.  Reconstruction gives
the first equality in (OBT.6).  Each $t_j$ is itself a
peripheral loop of the original surface subgroup $K$,
so the other four conjugacies are the actual peripheral
covariance of A.4.40.

There is also an exact gate for the remaining inner ambiguity.
Choose each ordinary centralizer representative on the branch
disk to fix $b_0$ as an element, as in the proof of (FBO.27).
If $a$ is such a representative, then
$\gamma a\gamma^{-1}$ fixes $b_0^2$, because
$\gamma(b_0^2)=b_0^{2\lambda}$ and $\lambda$ is a unit.
An ordinary representative of its transported outer class
also fixes $b_0$, and therefore $b_0^2$.
The two based orbifold representatives differ by
$\operatorname{Inn}(z)$ with $z\in P$; their inner
discrepancy consequently centralizes $b_0^2$.
Equations (OBT.1)--(OBT.2) imply that its
image in $Q_t$ belongs to $\langle b_0\rangle$, and is
therefore even for $\chi_T$.
Consequently restriction of that square to $K_*$ and then
to $W_s$ changes it only by an inner automorphism of
the actual surface quotient $W_s$.
This is the required outer covariance after the new
ordered-pole cover; it does not discard an odd orbifold
conjugator.

The finite cone lemma is proved independently in
\path{work/ghost_finite_cone_fusion.tex}.
The remaining assembly into the common translation-normalizer,
centralizer-recognition, and five-point-input level is recorded
separately.  The lemmas here are finite group statements with
their actual peripheral maps; no finite pentagon locus is
asserted to be a group.
